\chapter*{Foreword to the Second Printing of the Second Edition}
\addcontentsline{toc}{chapter}{Foreword}

This copy of the second edition of Seed Germination Theory and Practice is from
a second printing which accounts for any delays. The experiments on seed
germination is continuing and that has to take first priority. As compensation
you will get the following brief update on some of the more significant
observations made since the first printing in May 1993. Also the opportunity
was taken to correct a few errors. The following update briefly summarizes
some of the more significant results obtained since the first printing.

Alangium platanifolia seeds were washed and cleaned for seven days. Such seeds
germinated over several weeks at 70 when treated with GA-3 and not otherwise.
This is the first example of seeds in a fruit that had a GA-3 requirement for
germination.

Adonis vernalis and its geographical variants continue to be a problem. Large
amounts of seeds from two commercial houses have proven to be over 99\% empty
seed coats the same as seed from our own plantings and various other sources.
There is a need for someone to breed seed viability back into this species.
Several more cacti have been found to require GA-3 for germination. It is
possible and even likely that the majority of cacti will be found to have this
requirement. A definitive article has been written by myself and will appear in the
Journal of the Cactus and Succulent Society early in 1994.

Caltha palustris required GA-3 for germination. Seeds rapidly die when held
moist at 70, and all are dead in one month. This is one of the most rapid death rates
yet observed for seeds held moist at 70.

Caulophyllum thalictroides and Sanguinaria canadensis self sow on our property
yet fail to germinate under all conditions including treatment with GA-3. It is
likely that these species have a gibberelin requirement for germination, but
that some gibberelin(s) other than GA-3 are required. Note that it was already
shown that different gibberelins initiate germination in Shortia galacifolia
but in different patterns. The whole problem of gibberelin requirements for
germination are obviously complex. It will be some time before they are
completely understood.

Related to the gibberelin situation is the problem of germination in Trillium.
Trillium apetalon and Trillium pusillum have now been shown to require GA-3 for the
first step in germination, namely formation of the corm and root. Several other Trillium
had this behavior, but others had more complex responses to GA-3. Needliess to say
the Trillium problem is under intensive study.

Plentiful supplies of seeds of Lysochiton americanum and Symplocarpus
foetidus have now been obtained from our own plantings and the proper experiments
were conducted. Both were found to require light for germination typical of most
swamp plants. Germination occurs in several weeks in light at 70 and none in dark.

Seeds of four arctic willows (Salix arctica, S. herbaceae, S. lanata, and S;
phylicifolia) were collected by D. Haraldsson in Iceland and sent immediately. These
germinated in two days at 70. Dry storage for two weeks at 70 completely killed the
seed. However, the seeds can be stored moist at 40 for at least four weeks. Such
seeds germinate in two days the same as fresh seed when shifted to 70. It is of
interest that the seeds turn green when moist at 40, but do not develop further.
Sambucus pubens has been found to require GA-3 for germination the same as
other Sambucus. This seems to be a general requirement in this genus.
Senecio aureus required light for germination. Further, exposure to moist
conditions at 70 for just two months led to complete death of the seeds. Vernonia
altissima also required light (or GA-3) for germination at 70. It is unusual for such light
requirements in species of Asteraceae.

I am particularly indebted to Thomas Fischer, Senior Editor of Horticulture,
for a charmingly written account of my work. As a result the first printing of
the second edition sold out. Purchasers who are receiving the second printing
will forgive a somewhat delayed shipment. I also wish to express my
appreciation to the many amateur gardeners, seedsmen, and horticulturists who
have cooperated by sending me seeds from every corner of the World. Further,
it has been the utmost pleasure to see how nurseries and professional
propagators have embraced my book with enthusiasm. Many kind letters were
received for which I am most grateful. In contrast, there has been little
reaction from the academic community despite the fact that my work on seed
germination applies techniques more sophisticated and analysis more profound
than has been used in past work. There are reasons for this lack of response.
The problem faced by departments of horticulture and other biologists is that
they can no longer limit-their research to bending over microscopes.

To survive they are increasingly becoming pseudo-chemistry departments, because
this is the direction that has the potential for large research grants. The
previous head of the Department of Horticulture at Penn State had his Ph. D. in
chemistry. The Department of Botany, Microbiology, and Zoology are gone. These
are the signs of the times and indicative of the directions things are going.
The type of work described in this book is not the type that brings in big
research money despite its importance to both theory and practice of
horticulture. It will be interesting to see whether it ever gets much reaction
from the academic community.

Finally to end on a note of humility, on recounting the number of species
studied, the number is about 2500, not 4000 as given in the first printing. Further,
gibberellin and gibberellic acid were misspelled throughout and should have an added
letter I.
